\subsection{The theorem of the primitive element}

Let E be an extension of K; an element x of E is said to be primitive if \(E = K[x]\) . For the extension E to possess a primitive element it is necessary for \([E : K]\) to be finite.

THEOREM 1. --- Let E be an extension of K. Then the following conditions are equivalent :

\begin{enumerate}
\def\labelenumi{\alph{enumi})}
\item
  E possesses a primitive element;
\item
  there exist only a finite number of subextensions of E.
\end{enumerate}

These conditions are satisfied when E is a separable extension of finite degree.

Suppose first that E possesses a primitive element x, and let f be the minimal polynomial of x over K. For each monic polynomial \(g \in E[X]\) dividing f in E{[}X{]} denote by \(E_g\) the subextension of E generated by the coefficients of g. Since the possible polynomials g are finite in number (if f splits in E{[}X{]} into a product of r monic irreducible polynomials, the number is bounded by \(2^{r}\) ), the subextensions \(E_g\) are finite in number. To prove b) it is therefore enough to show that every subextension L of E is one of the \(E_g\) . Now if L is a subextension of E, then we have L{[}x{]} = E; if g is the minimal polynomial of x over L, we have {[}E:L{]} = deg(g). Besides, g is a divisor of f in L{[}X{]}, hence in E{[}X{]}; so we have \(E_g \subset L\) and \(E = E_g[x]\) . Since \(g(x) = 0\) , we have \([E : E_g] \leq \deg(g)\) , therefore \([E : E_g] \leq [E : L]\) and so \(L = E_g\) as we wished to show.

Next we observe that Condition b) implies that the extension E is of finite degree: by Remark 2 of V, p.~18 it suffices to prove that it is algebraic; now if z is an element of E transcendental over K then the subextensions \(\mathbf{K}(z^{n})\) , \(n \in N\) are pairwise distinct.

To show that \(b) \Rightarrow a)\) we now distinguish two cases:

\begin{enumerate}
\def\labelenumi{\Alph{enumi})}
\item
  If the field K is finite, the field E is a vector space of finite dimension over K and hence is a finite set. Therefore \({}^{1}\) (V, p.~78, Lemma 1) there exists an element x of E generating the multiplicative group of E, and we have \(E = K[x]\) .
\item
  Suppose now that the field K is infinite. If b) holds, the extension E is of finite degree, so b) can also be expressed by saying that E possesses only a finite number of subalgebras. This being so, the implication \(b) \Rightarrow a)\) is a consequence of the following more general proposition (for which the hypothesis that the field K be infinite is indispensable, cf.~V, p.~153, Ex. 5 of §7):
\end{enumerate}

PROPOSITION 7. --- Suppose that K is infinite; let A be a commutative K-algebra possessing only a finite number of sub-algebras (for example an etale K-algebra, V, p.~30, Prop. 3) and let V be a vector subspace of A generating A. Then there exists \(x \in V\) such that \(A = K[x]\) .

Let \(A_{1}, \ldots, A_{n}\) be the subalgebras of A distinct from A. If \(x \notin A_{1} \cup \ldots \cup A_{n}\) , then the sub-algebra \(K[x]\) cannot equal any of the \(A_{i}\) and so must coincide with A. Further, since V generates A, it is not contained in any of the subspaces \(A_{i}\) . Prop. 7 is therefore a consequence of the following lemma:

Lemma 1. --- Let A be a vector K-space, V, \(\mathbf{A}_1, \ldots, \mathbf{A}_n\) subspaces of A. If \(\operatorname{Card}(\mathbf{K}) \geqslant n\) and if V is not contained in any of the \(\mathbf{A}_i\) , then V is not contained in \(\mathbf{A}_1 \cup \ldots \cup \mathbf{A}_n\) .

Arguing by induction on \(n\) , we need only prove that if \(\mathbf{V} \subset \mathbf{A}_n\) and \(\mathbf{V} \subset \mathbf{A}_1 \cup \ldots \cup \mathbf{A}_n\) , then \(\mathbf{V} \subset \mathbf{A}_1 \cup \ldots \cup \mathbf{A}_{n-1}\) . Let \(x \in \mathbf{V}\) , \(x \notin \mathbf{A}_n\) , and let \(y\) be arbitrary in \(\mathbf{V}\) . If \(y \in \mathbf{K}x\) , we have \(y \in A_1 \cup \ldots \cup A_{n-1}\) ; if not, then the elements \(x\) and \(y + \lambda x\) , \(\lambda \in \mathbf{K}\) are strictly greater than \(n\) in number and belong to \(\mathbf{A}_1 \cup \ldots \cup \mathbf{A}_n\) , so two of them belong to the same \(\mathbf{A}_i\) . Thus there exists \(i\) , \(1 \leqslant i \leqslant n\) such that either \(x \in \mathbf{A}_i\) and \(y + \lambda x \in \mathbf{A}_i\) for some \(\lambda \in \mathbf{K}\) , or \(y + \lambda x \in \mathbf{A}_i\) and \(y + \mu x \in \mathbf{A}_i\) for two distinct scalars \(\lambda, \mu \in \mathbf{K}\) . In both cases we conclude that \(x \in \mathbf{A}_i\) and \(y \in \mathbf{A}_i\) ; but this implies that \(i \neq n\) , hence \(y \in \mathbf{A}_1 \cup \ldots \cup A_{n-1}\) as we wished to show.

This completes the proof of the equivalence of \(a)\) and \(b)\) in Th. 1. Finally if the extension \(E\) is separable and of finite degree, condition \(b)\) holds, by \(V\) , p.~30, Prop. 3.

